% 解答生成済。問題本体は大問ごとに1ページ。


% レビュー反映：疑問文の書写・和訳には、主語に対応する短い返答を添える。
\MajorQuestionLesson{三人称単数現在の否定文・疑問文}
\TypeQuestion{英文や会話文を書き写し、日本語に訳しなさい。}{generalverb_third_question_negative_copy_translate}
\SetRowSpace{3mm}
\EngCopyTranslateSplit{He doesn't play tennis.}{彼はテニスをしません。}
\EngCopyTranslateSplit{She does not like music.}{彼女は音楽が好きではありません。}
\EngCopyTranslate{Does Ken study English?\enspace---\enspace Yes, he does.}{ケンは英語を勉強しますか。\par はい、勉強します。}
\EngCopyTranslate{Does she have a pen?\enspace---\enspace No, she doesn't.}{彼女はペンを持っていますか。\par いいえ、持っていません。}
\EngCopyTranslate{Do they play tennis?\enspace---\enspace Yes, they do.}{彼らはテニスをしますか。\par はい、します。}

\LessonPageBreak
\TypeQuestion{英文を書き写したあと、主語に下線を引き、動詞を丸で囲みなさい。\par{\small ※会話文では、両方の文に印を付けること。}}{generalverb_third_question_negative_copy_structure}
\SetRowSpace{2.5mm}
\EngWriteQuestion{\EngDialogue{Does your sister play tennis?}{Yes, she does.}}{\EngAnswerPair{Does \underline{your sister} \EngCircle{play} tennis?}{\EngText{Yes, \underline{she} \EngCircle{does}.}}}
\EngWriteQuestion{\EngSentence{He doesn't study English.}}{\underline{He} doesn't \EngCircle{study} English.}
\EngWriteQuestion{\EngDialogue{Are they busy?}{No, they aren't.}}{\EngAnswerPair{\EngCircle{Are} \underline{they} busy?}{\EngText{No, \underline{they} \EngCircle{aren't}.}}}
\EngWriteQuestion{\EngDialogue{Does Ken have a book?}{No, he doesn't.}}{\EngAnswerPair{Does \underline{Ken} \EngCircle{have} a book?}{\EngText{No, \underline{he} \EngCircle{doesn't}.}}}
\EngWriteQuestion{\EngSentence{We don't watch TV.}}{\underline{We} don't \EngCircle{watch} TV.}
\EngWriteQuestion{\EngSentence{She goes to school.}}{\underline{She} \EngCircle{goes} to school.}
\EngWriteQuestion{\EngDialogue{What does your brother study?}{He studies English.}}{\EngAnswerPair{What does \underline{your brother} \EngCircle{study}?}{\EngText{\underline{He} \EngCircle{studies} English.}}}
\EngWriteQuestion{\EngSentence{My sisters are kind.}}{\underline{My sisters} \EngCircle{are} kind.}
\EngWriteQuestion{\EngSentence{She doesn't do her homework.}}{\underline{She} doesn't \EngCircle{do} her homework.}
\EngWriteQuestion{\EngDialogue{Do you like music?}{Yes, I do.}}{\EngAnswerPair{Do \underline{you} \EngCircle{like} music?}{\EngText{Yes, \underline{I} \EngCircle{do}.}}}

\LessonPageBreak
\TypeQuestion{日本語に合うように、誤りのある英文の必要なところだけ直して書きなさい。\par{\small ※短縮形の有無は原文に合わせること。右の返答は参考にすること。}}{generalverb_third_question_negative_error_correct}
\SetRowSpace{2mm}
\EngWriteQuestion{\EngDialogue{Does he plays tennis?}{Yes, he does.}\par〔彼はテニスをしますか。\enspace---\enspace はい、します。〕}{Does he play tennis?}
\EngWriteQuestion{\EngSentence{She don't like music.}\par〔彼女は音楽が好きではありません。〕}{She doesn't like music.}
\EngWriteQuestion{\EngSentence{He doesn't has a pen.}\par〔彼はペンを持っていません。〕}{He doesn't have a pen.}
\EngWriteQuestion{\EngDialogue{Is your sister play tennis?}{Yes, she does.}\par〔あなたの姉はテニスをしますか。\enspace---\enspace はい、します。〕}{Does your sister play tennis?}
\EngWriteQuestion{\EngSentence{They doesn't watch TV.}\par〔彼らはテレビを見ません。〕}{They don't watch TV.}
\EngWriteQuestion{\EngDialogue{What does she studies?}{She studies English.}\par〔彼女は何を勉強しますか。\enspace---\enspace 彼女は英語を勉強します。〕}{What does she study?}
\EngWriteQuestion{\EngSentence{He does not busy.}\par〔彼は忙しくありません。〕}{He is not busy.}
\EngWriteQuestion{\EngDialogue{Does she does her homework?}{No, she doesn't.}\par〔彼女は宿題をしますか。\enspace---\enspace いいえ、しません。〕}{Does she do her homework?}

\LessonPageBreak
\TypeQuestion{日本語に合うように、適切な語句を丸で囲み、英文を書きなさい。\par{\small ※会話文では、疑問文と返答の両方を完成させて書くこと。}}{generalverb_third_question_negative_fill_aux_choice}
\SetRowSpace{1.5mm}
\EngWriteQuestion{彼はテニスをしません。\par\EngSentence{He \EngChoice{doesn't / don't} play tennis.}}{He doesn't play tennis.}
\EngWriteQuestion{あなたの姉は英語を勉強しますか。\enspace---\enspace はい、勉強します。\par\EngDialogue{\EngChoice{Do / Does} your sister study English?}{Yes, she \EngChoice{do / does}.}}{\EngAnswerPair{Does your sister study English?}{\EngText{Yes, she does.}}}
\EngWriteQuestion{彼らは音楽が好きではありません。\par\EngSentence{They \EngChoice{don't / doesn't} like music.}}{They don't like music.}
\EngWriteQuestion{彼女は忙しくありません。\par\EngSentence{She \EngChoice{isn't / doesn't} busy.}}{She isn't busy.}
\EngWriteQuestion{ケンは本を持っていますか。\enspace---\enspace いいえ、持っていません。\par\EngDialogue{\EngChoice{Is / Does} Ken have a book?}{No, he \EngChoice{isn't / doesn't}.}}{\EngAnswerPair{Does Ken have a book?}{\EngText{No, he doesn't.}}}
\EngWriteQuestion{私たちはテレビを見ません。\par\EngSentence{We \EngChoice{aren't / don't} watch TV.}}{We don't watch TV.}
\EngWriteQuestion{あなたの兄は親切ですか。\enspace---\enspace はい、親切です。\par\EngDialogue{\EngChoice{Does / Is} your brother kind?}{Yes, he \EngChoice{does / is}.}}{\EngAnswerPair{Is your brother kind?}{\EngText{Yes, he is.}}}
\EngWriteQuestion{ユイは宿題をしません。\par\EngSentence{Yui \EngChoice{doesn't / isn't} do her homework.}}{Yui doesn't do her homework.}

\LessonPageBreak
\TypeQuestion{日本語に合うように、適切な語句を丸で囲み、英文を書きなさい。\par{\small ※会話文では、疑問文と返答の両方を完成させて書くこと。}}{generalverb_third_question_negative_fill_verb_choice}
\SetRowSpace{1.5mm}
\EngWriteQuestion{彼女は英語を勉強しますか。\enspace---\enspace いいえ、勉強しません。\par\EngDialogue{Does she \EngChoice{studies / study} English?}{No, she \EngChoice{don't / doesn't}.}}{\EngAnswerPair{Does she study English?}{\EngText{No, she doesn't.}}}
\EngWriteQuestion{彼は学校に行きません。\par\EngSentence{He doesn't \EngChoice{go / goes} to school.}}{He doesn't go to school.}
\EngWriteQuestion{ケンはテレビを見ます。\par\EngSentence{Ken \EngChoice{watch / watches} TV.}}{Ken watches TV.}
\EngWriteQuestion{あなたの姉はペンを持っていますか。\enspace---\enspace はい、持っています。\par\EngDialogue{Does your sister \EngChoice{has / have} a pen?}{Yes, she \EngChoice{do / does}.}}{\EngAnswerPair{Does your sister have a pen?}{\EngText{Yes, she does.}}}
\EngWriteQuestion{彼女は音楽が好きではありません。\par\EngSentence{She doesn't \EngChoice{likes / like} music.}}{She doesn't like music.}
\EngWriteQuestion{彼らは毎日テニスをします。\par\EngSentence{They \EngChoice{plays / play} tennis every day.}}{They play tennis every day.}
\EngWriteQuestion{彼は宿題をしますか。\enspace---\enspace いいえ、しません。\par\EngDialogue{Does he \EngChoice{do / does} his homework?}{No, he \EngChoice{doesn't / don't}.}}{\EngAnswerPair{Does he do his homework?}{\EngText{No, he doesn't.}}}
\EngWriteQuestion{私の姉は忙しくありません。\par\EngSentence{My sister \EngChoice{does / is} not busy.}}{My sister is not busy.}

\LessonPageBreak
\TypeQuestion{疑問文に対する答えとして最も適切なものを一つ選び、その記号を丸で囲みなさい。\par 疑問文と選んだ返答を、続けて日本語に訳しなさい。}{generalverb_third_question_negative_response_choice}
\SetRowSpace{2mm}
\EngResponseChoice{Does he play tennis?}{\mbox{（ア）\EngText{Yes, he do.}}\qquad\mbox{（イ）\EngText{Yes, he does.}}\qquad\mbox{（ウ）\EngText{Yes, he is.}}}{\parbox[t]{\linewidth}{イ\par 彼はテニスをしますか。\par はい、します。}}
\EngResponseChoice{Does she like music?}{\mbox{（ア）\EngText{No, she don't.}}\qquad\mbox{（イ）\EngText{No, he doesn't.}}\qquad\mbox{（ウ）\EngText{No, she doesn't.}}}{\parbox[t]{\linewidth}{ウ\par 彼女は音楽が好きですか。\par いいえ、好きではありません。}}
\EngResponseChoice{Do they study English?}{\mbox{（ア）\EngText{Yes, they do.}}\qquad\mbox{（イ）\EngText{Yes, they does.}}\qquad\mbox{（ウ）\EngText{Yes, they are.}}}{\parbox[t]{\linewidth}{ア\par 彼らは英語を勉強しますか。\par はい、勉強します。}}
\EngResponseChoice{Is your sister busy?}{\mbox{（ア）\EngText{Yes, she does.}}\qquad\mbox{（イ）\EngText{Yes, you are.}}\qquad\mbox{（ウ）\EngText{Yes, she is.}}}{\parbox[t]{\linewidth}{ウ\par あなたの姉は忙しいですか。\par はい、忙しいです。}}
\EngResponseChoice{Does Ken have a pen?}{\mbox{（ア）\EngText{No, he isn't.}}\qquad\mbox{（イ）\EngText{No, he doesn't.}}\qquad\mbox{（ウ）\EngText{No, I don't.}}}{\parbox[t]{\linewidth}{イ\par ケンはペンを持っていますか。\par いいえ、持っていません。}}
\EngResponseChoice{Does your brother watch TV?}{\mbox{（ア）\EngText{No, he doesn't.}}\qquad\mbox{（イ）\EngText{No, he don't.}}\qquad\mbox{（ウ）\EngText{No, they don't.}}}{\parbox[t]{\linewidth}{ア\par あなたの兄はテレビを見ますか。\par いいえ、見ません。}}
\EngResponseChoice{Does she do her homework?}{\mbox{（ア）\EngText{Yes, she is.}}\qquad\mbox{（イ）\EngText{Yes, she does.}}\qquad\mbox{（ウ）\EngText{Yes, I do.}}}{\parbox[t]{\linewidth}{イ\par 彼女は宿題をしますか。\par はい、します。}}
\EngResponseChoice{What does Ken study?}{\mbox{（ア）\EngText{English.}}\qquad\mbox{（イ）\EngText{At school.}}\qquad\mbox{（ウ）\EngText{At seven.}}}{\parbox[t]{\linewidth}{ア\par ケンは何を勉強しますか。\par 英語です。}}

\LessonPageBreak


% レビュー反映：三単現を問う「すべて使用・1語不要・1語不足」とandを並べる問題を混在。
\TypeQuestion{日本語に合う英文・会話文を、語句を並べ替えて作りなさい。\par{\small ※返答の語群はすべて使用。人名は日本語順、時の語句は文末に。}}{generalverb_third_question_negative_reorder}
\SetRowSpace{1.5mm}
\EngWriteQuestion{彼女はテニスをしません。\quad〔1語不要〕\par\EngSentence{\EngReorder{play / don't / she / tennis / doesn't}}}{She doesn't play tennis.}
\EngWriteQuestion{ケンは英語を勉強しますか。\enspace---\enspace はい、勉強します。\quad〔1語不足〕\par\EngDialogue{\EngReorder{study / Ken / English}}{\EngReorder{he / does / yes}}}{\EngAnswerPair{Does Ken study English?}{\EngText{Yes, he does.}}}
\EngWriteQuestion{彼はペンを持っていません。\quad〔すべて使用〕\par\EngSentence{\EngReorder{have / he / a pen / does / not}}}{He does not have a pen.}
\EngWriteQuestion{あなたの姉はテレビを見ますか。\enspace---\enspace いいえ、見ません。\quad〔1語不要〕\par\EngDialogue{\EngReorder{TV / your sister / watch / is / does}}{\EngReorder{she / no / doesn't}}}{\EngAnswerPair{Does your sister watch TV?}{\EngText{No, she doesn't.}}}
\EngWriteQuestion{彼女は何を読みますか。\enspace---\enspace 彼女は本を読みます。\quad〔すべて使用〕\par\EngDialogue{\EngReorder{read / she / what / does}}{\EngReorder{books / reads / she}}}{\EngAnswerPair{What does she read?}{\EngText{She reads books.}}}
\EngWriteQuestion{ケンとユイは英語を勉強しますか。\enspace---\enspace はい、勉強します。\quad〔1語不足〕\par\EngDialogue{\EngReorder{study / Yui / Ken / English / and}}{\EngReorder{do / yes / they}}}{\EngAnswerPair{Do Ken and Yui study English?}{\EngText{Yes, they do.}}}
\EngWriteQuestion{彼は親切ではありません。\quad〔すべて使用〕\par\EngSentence{\EngReorder{not / he / kind / is}}}{He is not kind.}
\EngWriteQuestion{ユイは宿題をしません。\quad〔1語不足〕\par\EngSentence{\EngReorder{her homework / Yui / not / do}}}{Yui does not do her homework.}

\LessonPageBreak
\TypeQuestion{英文の右の指示に従って、文を書き換えなさい。}{generalverb_third_question_negative_transform}
\SetRowSpace{3mm}
\EngWriteQuestion{\EngSentence{He plays tennis.}\quad〔短縮形を使わない否定文に〕}{He does not play tennis.}
\EngWriteQuestion{\EngSentence{She studies English.}\quad〔疑問文にし、\EngText{Yes}で短く答える〕}{\EngAnswerPair{Does she study English?}{\EngText{Yes, she does.}}}
\EngWriteQuestion{\EngSentence{Ken has a book.}\quad〔短縮形を使った否定文に〕}{Ken doesn't have a book.}
\EngWriteQuestion{\EngSentence{Your sister watches TV.}\quad〔疑問文にし、\EngText{No}で短く答える〕}{\EngAnswerPair{Does your sister watch TV?}{\EngText{No, she doesn't.}}}
\EngWriteQuestion{\EngSentence{He doesn't like music.}\quad〔肯定文に〕}{He likes music.}
\EngWriteQuestion{\EngSentence{They play tennis.}\quad〔疑問文にし、\EngText{Yes}で短く答える〕}{\EngAnswerPair{Do they play tennis?}{\EngText{Yes, they do.}}}
\EngWriteQuestion{\EngSentence{She is busy.}\quad〔短縮形を使わない否定文に〕}{She is not busy.}
\EngWriteQuestion{\EngSentence{My brother goes to school.}\quad〔疑問文にし、\EngText{No}で短く答える〕}{\EngAnswerPair{Does my brother go to school?}{\EngText{No, he doesn't.}}}
\EngWriteQuestion{\EngSentence{She does her homework.}\quad〔短縮形を使った否定文に〕}{She doesn't do her homework.}
\EngWriteQuestion{\EngSentence{He does not read books.}\quad〔短縮形を使って〕}{He doesn't read books.}

\WordBank{三人称単数現在の否定文・疑問文}{
  \WB{I / you}{代名詞}{私は／あなた（たち）は}{\EngText{I}は必ず大文字}
  \WB{we / they}{代名詞}{私たちは／彼らは}{\EngText{they}は物の複数にも使う}
  \WB{he / she / it}{代名詞}{彼は／彼女は／それは}{}
  \WB{my / your}{所有格}{私の／あなた（たち）の}{}
  \WB{his / her}{所有格}{彼の／彼女の}{}
  \WB{a}{冠詞}{1つの／1人の}{訳さない場合もある}
  \WB{am / are / is}{be動詞}{～です／～にいます}{主語に合わせて使い分ける}
  \WB{do}{動詞等}{する；疑問・否定を作る}{三単現 \EngText{does}}
  \WB{does}{動詞等}{する；疑問・否定を作る}{\EngText{do}の三単現}
  \WB{not}{副詞}{～でない／～しない}{}
  \WB{don't}{短縮形}{しない；いいえ、しない}{\EngText{do not}の短縮形}
  \WB{doesn't}{短縮形}{～しない}{\EngText{does not}の短縮形}
  \WB{aren't}{短縮形}{～でない}{\EngText{are not}の短縮形}
  \WB{isn't}{短縮形}{～でない}{\EngText{is not}の短縮形}
  \WB{yes / no}{応答語}{はい／いいえ}{}
  \WB{what}{疑問詞}{何；何の}{}
  \WB{play}{動詞}{（スポーツを）する}{三単現 \EngText{plays}}
  \WB{study}{動詞}{勉強する}{三単現 \EngText{studies}}
  \WB{like}{動詞}{～が好きである}{三単現 \EngText{likes}}
  \WB{read}{動詞}{読む}{三単現 \EngText{reads}}
  \WB{have}{動詞}{持っている；食べる}{三単現 \EngText{has}}
  \WB{go}{動詞}{行く}{三単現 \EngText{goes}}
  \WB{watch}{動詞}{見る}{三単現 \EngText{watches}}
  \WB{Ken / Yui}{人名}{ケン／ユイ}{最初の文字は大文字}
  \WB{brother}{名詞}{兄；弟}{複数形 \EngText{brothers}}
  \WB{sister}{名詞}{姉；妹}{複数形 \EngText{sisters}}
  \WB{friend}{名詞}{友達}{複数形 \EngText{friends}}
  \WB{student}{名詞}{生徒}{複数形 \EngText{students}}
  \WB{English}{名詞}{英語}{最初の文字は大文字}
  \WB{music}{名詞}{音楽}{}
  \WB{tennis}{名詞}{テニス}{}
  \WB{book}{名詞}{本}{複数形 \EngText{books}}
  \WB{pen}{名詞}{ペン}{}
  \WB{school}{名詞}{学校}{}
  \WB{homework}{名詞}{宿題}{}
  \WB{TV}{名詞}{テレビ}{\EngText{television}の略}
  \WB{day}{名詞}{日；1日}{}
  \WB{busy}{形容詞}{忙しい}{}
  \WB{kind}{形容詞}{親切な}{}
  \WB{every}{形容詞}{毎～}{}
  \WB{seven}{数詞}{7（七）}{}
  \WB{at}{前置詞}{～で／～時に}{}
  \WB{to}{前置詞}{～へ／～に}{}
  \WB{and}{接続詞}{～と…}{}
  \WB{every day}{熟語}{毎日}{}
}

\PhraseBank{三人称単数現在の否定文・疑問文}{
  \PB{He does not play tennis.}{彼はテニスをしません。}
  \PB{He doesn't play tennis.}{彼はテニスをしません。}
  \PB{She doesn't have a pen.}{彼女はペンを持っていません。}
  \PB*{Does she study English?}{彼女は英語を勉強しますか。}
  \PB*{Yes, she does.}{はい、勉強します。}
  \PB{No, she doesn't.}{いいえ、勉強しません。}
  \PB*{Does Ken have a book?}{ケンは本を持っていますか。}
  \PB{Yes, he does.}{はい、持っています。}
  \PB*{What does he study?}{彼は何を勉強しますか。}
  \PB{He studies English.}{彼は英語を勉強します。}
  \PB*{Does she do her homework?}{彼女は宿題をしますか。}
  \PB{No, she doesn't.}{いいえ、しません。}
}
